% Perry Anderson: Ever Closer Union?: Reading Summary
% Introduction to the European Union, Week 4
% Compile with: pdflatex anderson-ever-closer-union-reading-summary.tex

\documentclass[11pt,a4paper]{article}
\usepackage[utf8]{inputenc}
\usepackage[T1]{fontenc}
\usepackage[margin=2.5cm]{geometry}
\usepackage{booktabs}
\usepackage{enumitem}
\setlist{nosep,leftmargin=*}
\usepackage{parskip}
\usepackage{microtype}
\usepackage{hyperref}
\usepackage{xcolor}
\definecolor{eublue}{RGB}{0,51,153}

\hyphenation{Spitzenkandidat Spitzenkandidaten}
\hypersetup{
  colorlinks=true,
  linkcolor=eublue,
  urlcolor=eublue,
  pdftitle={Anderson: Ever Closer Union?},
  pdfauthor={Introduction to the EU, UCM}
}

\begin{document}

\title{Perry Anderson: Ever Closer Union?\\Critical Analysis of European Integration}
\author{Introduction to the European Union\\Universidad Complutense de Madrid\\Week 4}
\date{}
\maketitle
\vspace{-1em}
\hrule
\vspace{1em}

\section*{Rhetorical Framework (Ethos, Logos, Pathos)}

\begin{table}[h]
\centering
\begin{tabular}{@{}lll@{}}
\toprule
\textbf{Appeal} & \textbf{Meaning} & \textbf{Anderson's Use} \\
\midrule
\textbf{Ethos} & Credibility, authority, trust & Left-wing intellectual; critical theory; contrarian voice against celebratory narratives \\
\textbf{Logos} & Logic, structure, argument & Systematic critique; identification of tensions, contradictions, uneven integration \\
\textbf{Pathos} & Emotion, stakes, persuasion & Disillusionment, urgency; appeal to readers skeptical of official EU narratives \\
\bottomrule
\end{tabular}
\end{table}

\textbf{Key insight:} Anderson deliberately challenges \textit{ethos} of pro-European scholarship---he positions himself as critical outsider. His \textit{logos} (tensions, unevenness, limits) supports that stance. \textit{Pathos} targets readers who share skepticism about ``ever closer union'' or feel the EU has failed to deliver.

\section*{Bibliographic Information}

\begin{itemize}
\item \textbf{Author:} Perry Anderson
\item \textbf{Title:} Ever Closer Union?
\item \textbf{Year:} 2021
\end{itemize}

\section*{Summary \& Key Points}

\begin{itemize}
\item Anderson provides a critical analysis of European integration, questioning the ``ever closer union'' narrative.
\item Examines the tensions, contradictions, and limitations of the European project.
\item Analyzes how integration has proceeded unevenly across different policy areas.
\item Questions whether ``ever closer union'' is still a viable or desirable goal.
\item Provides critical perspective on current challenges facing the European Union.
\item Connects historical development to contemporary debates about European integration.
\end{itemize}

\textbf{Rhetorical note:} The title itself---\textit{Ever Closer Union?}---is interrogative: \textit{logos} (question as argument) + \textit{pathos} (doubt, skepticism). Anderson inverts the official motto into a question.

\section*{Main Arguments}

\subsection*{I. Critical Assessment of Integration}

Anderson questions the assumption that European integration necessarily leads to ``ever closer union.'' The title itself is interrogative, suggesting doubt about this trajectory.

\subsection*{II. Uneven Integration}

The reading examines how integration has proceeded unevenly---some areas (like economic integration) have advanced further than others (like political union or social policy).

\subsection*{III. Contemporary Challenges}

As a 2021 publication, Anderson's work addresses current challenges facing the EU, including Brexit, democratic deficits, and tensions between member states.

\subsection*{IV. Critical Perspective}

Anderson provides a critical, left-wing perspective on European integration, questioning both the process and outcomes of integration from a critical theory standpoint.

\section*{Context \& Analysis}

\textbf{Strengths:} Provides critical perspective; addresses contemporary challenges; questions assumptions about integration. Complements Krawatzek \& Pestel's call for diverse perspectives.

\textbf{Weaknesses:} May be overly critical or pessimistic; may not fully acknowledge achievements of integration. Risk of \textit{pathos} overwhelming \textit{logos}.

\textbf{Connections:} Relates to critical historiography (Krawatzek \& Pestel); complements Vinen's fragmented approach; contrasts with Judt's pragmatic narrative; with Jarausch's defensive account.

\section*{Relevance to Course}

\begin{itemize}
\item \textbf{Ethos:} Essential for understanding critical perspectives; counterweight to celebratory narratives.
\item \textbf{Logos:} Tensions, contradictions, uneven integration---analytical framework.
\item \textbf{Pathos:} Stakes of contemporary debates; disillusionment, urgency.
\item \textbf{All three:} Course theme of presenting both official and critical perspectives; Anderson embodies the critical side.
\end{itemize}

\section*{Discussion Questions}

\begin{itemize}
\item What does Anderson mean by questioning ``Ever Closer Union?'' (Ethos: his position vs.\ official; Logos: main critiques; Pathos: tone, stakes)
\item How has integration proceeded unevenly? (Ethos: whose perspective? Logos: evidence, policy areas; Pathos: what's missing?)
\item Main contemporary challenges? (Ethos: Anderson's diagnosis; Logos: evidence, causes; Pathos: urgency)
\item Is ``ever closer union'' viable or desirable? (Ethos: who decides? Logos: arguments for/against; Pathos: hope vs.\ doubt)
\item How does Anderson differ from celebratory narratives? (Ethos: voice, stance; Logos: contrasting claims; Pathos: emotional register)
\end{itemize}

\vspace{1em}
\noindent\footnotesize\textit{Reference:} Anderson, Perry. \textit{Ever Closer Union?} 2021.

\end{document}
